\section{Literature review}
Despite its disadvantages, indirect methods are widely used when computing space flight trajectories. The main reason is the low accuracy associated with direct methods. In practice when using direct methods the minimum functional value found can be expected to have an error of about one percent which can be a crucial part of the payload in a space flight mission. Increasing the finite-dimensional space doesn't necessarily yield better results \cite{von1992direct}.

One such indirect method is the \acrfull{sgra}. Developed during the years 1968-86 by Miele \emph{et al.} \cite{miele1970sequential}, this method has been successfully applied to atmospheric flight \cite{miele1986optimal}, space flight \cite{miele1999optimaltransfers, miele2001optimaltrajectories} and flights at the interface between atmosphere and space \cite{miele1998optimal, miele1993nominal, miele2000feasibility, miele2001assessment}. Applications and extensions have been developed all around the world with a notable example being the implementation of \acrshort{sgra} used at the NASA Johnson Space Center under the code name SeGRAm \cite{rishikof1992segram, miele2007optimal}.

\acrshort{sgra} is a sequential algorithm that alternates between gradient and restoration phases. In the gradient phase variations on the decision variables are computed that reduce the value of the augmented objective function, while avoiding excessive constraint violations. Then in the restoration phase variations on the decision variables are computed that reduce the constraint errors while avoiding excessive change in the decision variables.

In 2003 Miele \emph{et al.} developed an extension to \acrshort{sgra} called \acrfull{msgra}. The reason was to improve the robustness of the method and enlarge the field of applications. \acrshort{msgra} can be applied to optimal control problems involving multiple subsystems as well as discontinuities in the state and control variables at the interface between contiguous subsystems \cite{miele2003multiplep1}. \acrshort{msgra} has demonstrated in being an effective tool for optimizing the trajectories of multiple stage launch vehicles \cite{miele2003multiplep2}.

A mayor disadvantage of indirect methods is the high sensitivity and difficulty to select an appropriate initial guess. In order to lessen this problem some applications apply the use of the homotopy method. The principle of the homotopy method is that a given problem is embedded into a family of problems parametrized by a homotopic parameter. The solution to the original problem can then be found by starting with a simpler problem and gradually adjust the homotopic parameter such that in the end a solution to the original problem is found \cite{pan2016double}. However the issue with the homotopy method is the added complexity to the solution process.

As an example Calise \emph{et al.} \cite{calise2000further} and, Gath and Calise \cite{gath2001optimization} propose the use of a hybrid analytical/numerical approach which uses a homotopy method that starts with a nearly analytic vacuum solution that gradually introduces atmospheric effects until a converged solution is obtained. This method has the potential to be used online, since once converged, the trajectory can be updated within milliseconds using input from the guidance system due to the nearly analytic nature of the solution process \cite{gath2001optimization}. 

Direct methods have also been proposed, such as those by Roh and Kim \cite{roh2002trajectory} and Jamilnia and Naghash \cite{jamilnia2012simultaneous}, both of whom solve the problem by using direct collocation to transcribe the trajectory optimization problem into \acrfull{nlp} problems. Another direct method proposed by Coverstone-Carrol and Williams \cite{coverstone1994optimal} uses differential inclusion concepts to remove explicit control dependence from the problem statement thereby reducing the number unknown parameters.

Recently newer techniques make use of heuristic optimization methods. A problem inherent to the methods already discussed is that since these methods are based on local computations, they are sensitive the starting guess and do not always converge towards the global optimum. On the other hand heuristic methods are global techniques which can look for the global optimum within a larger portion of the state space. The main idea behind heuristic optimization methods is that the search is done in a stochastic manner instead of a deterministic manner \cite{rao2009survey}. However the main disadvantage of heuristic methods is the lack of any convergence proof and the capability of yielding only near optimal solutions \cite{pontani2014optimal}.

These disadvantages can be circumvented by combining indirect and heuristic methods. As an example Pontani \emph{et al.} \cite{PONTANI2014852, pontani2013ascent, pontani2014simple, pontani2014optimal} propose a method for computing the optimal exoatmospheric trajectory using \acrfull{pso}. The proposed method starts by analytically deriving the necessary conditions for optimality just as with indirect methods, then \acrshort{pso} is used to numerically determine the state, control and adjoint variables. The use of a heuristic method significantly lessens the initialization issues while retaining the accuracy and measure of accuracy that come with indirect methods, without the added complexity associated with the use of homotopic methods. 

A more direct heuristic optimization approach is that taken by Zheng \emph{et al.} \cite{zheng2020ascent}. A pitch law is defined consisting of five successive sections, each for a single phase of flight with the first one starting at take-off and the last one ending at orbit insertion. A function with unknown parameters is chosen for each section of the pitch law, and then the unknowns are found using \acrfull{ga}. Since these functions can be selected based on prior knowledge, it's possible to significantly reduce the number of unknowns. In this case there where only nine.

This is a similar approach as to that taken by Pontani and Teofilatto \cite{pontani2014simple} to compute the endoatmospheric flight trajectory using orbital ascent. However in this case, the parameters where found using gradient based optimization methods instead of heuristic methods.



